% ----------------------------------------------------------------
\section{Guarantees and Correctness Properties}
\label{sec:correctness}
% ----------------------------------------------------------------

The preceding sections specify the Bailout Protocol as a deterministic state machine
$\mathcal{B}(n)$, a trigger predicate, an escalation algorithm, and a bypass mechanism.
This section states the three correctness properties the protocol must guarantee,
identifies the structural assumptions each rests on, and provides a proof sketch for
each. The properties are stated as invariants over the full reachable state space of
$\mathcal{B}$, not merely for a specific execution path. Together they constitute the
protocol's correctness claim: every unresolved exception in any \bxthree{} node either
resolves within the node's bounded authority or compulsorily transfers control to a
named human principal, and this transfer is structurally indestructible by any machine
actor.

% ----------------------------------------------------------------
\subsection{Structural Assumptions}
\label{sec:assumptions}

All three proofs rely on the following structural facts established in
Sections~\ref{sec:bailout} and~\ref{sec:bailout-trigger}:

\begin{enumerate}
  \setlength\itemsep{0pt}
  \item[\textbf{(A1)}] \label{asm:fact-layer} The \fact{} layer pre-commit hook is
    operative: it is the sole write authority for the Forensic Ledger and rejects any
    ledger update that would record a directive issued by a machine actor.
  \item[\textbf{(A2)}] \label{asm:immutable-pointer} The immutable parent pointer
    $\alpha(n)$ and human accountability anchor $h(n)$ are established at node
    provisioning and are read-only at the machine-actor level for the node's operational
    lifetime.
  \item[\textbf{(A3)}] \label{asm:determinism} The state machine $\mathcal{B}(n)$ is
    deterministic: for every $(q, \sigma) \in Q \times \Sigma$, at most one
    $q' \in Q$ satisfies $(q, \sigma, q') \in \rightarrow$.
  \item[\textbf{(A4)}] \label{asm:pathways} At least one functional communication
    pathway to $h(n)$ is provisioned at node startup. This is a deployment prerequisite,
    not a protocol invariant.
\end{enumerate}

These assumptions are guaranteed by the \bxthree{} framework's three-layer architecture
and by the provisioning process at node initialization. They do not depend on any
property of the machine actors that execute within the system.

% ----------------------------------------------------------------
\subsection{SAFETY: No Autonomous Collapse}
\label{sec:safety}

\textbf{Theorem (Safety).} \emph{No execution of the Bailout Protocol state machine
$\mathcal{B}$ can reach a terminal state in which an unresolved exception propagates
exclusively through machine actors, or in which an autonomous action with potential for
harm executes without explicit, recorded human authorization.}

\medskip
\noindent\textbf{Formal statement.} Let $\mathcal{B}$ be the state machine for node $n$.
Let $P(t)$ denote the set of functional communication pathways to $h(n)$ available at
time $t$. Safety requires:

\begin{equation}
\label{eq:safety}
\forall t : P(t) = \emptyset \;\Longrightarrow\;
\text{state}(n) \neq \text{RESOLVED}
\;\lor\;
\exists\, d \in \{\text{ACK},\ \text{RESOLVE},\ \text{REJECT}\} \;\land\; \text{issuer}(d) = h(n).
\end{equation}

Equivalently: the system cannot enter the \text{RESOLVED} state through a machine-issued
directive, regardless of pathway availability. If no pathway to the human anchor
exists, the system stalls in a non-terminal state and awaits human intervention or
pathway restoration.

\medskip
\noindent\textbf{Proof sketch.} The proof proceeds by contradiction over all reachable
states of $\mathcal{B}$.

\begin{itemize}
  \setlength\itemsep{0pt}
  \item[$\bullet$] \textbf{Step 1: The only terminal transitions.}
    The \text{RESOLVED} state is reachable only via T6
    ($\text{ESCALATING} \rightarrow \text{RESOLVED}$) or T7
    ($\text{HUMAN\_ACKNOWLEDGED} \rightarrow \text{RESOLVED}$). Both require a directive
    $d \in \{\text{RESOLVE}, \text{REJECT}\}$ issued by $h(n)$. By invariant
    \texttt{INV}$_3$ (Section~\ref{sec:bailout-trigger}), every \text{RESOLVED} state
    carries a recorded directive from the human anchor. A machine-issued directive
    cannot satisfy the preconditions for T6 or T7.

  \item[$\bullet$] \textbf{Step 2: Entry to \text{ESCALATING} is machine-actuator-resistant.}
    Transition T4 ($\text{BAIL\_PENDING} \rightarrow \text{ESCALATING}$) fires
    immediately upon entry to \text{BAIL\_PENDING} with \emph{no dwell time} and no
    machine-actor intervention permitted. This is an architectural compulsion, not a
    configurable delay. A machine actor cannot suppress T4 or absorb the exception at
    \text{BAIL\_PENDING}.

  \item[$\bullet$] \textbf{Step 3: The bypass mechanism blocks machine-termination.}
    The bypass property (BP) — proved in Section~\ref{sec:bypass} — guarantees that no
    machine actor along the parent chain may serve as the terminal resolver for an
    escalated exception. The escalation algorithm terminates exclusively at $h(n)$. A
    machine actor cannot route the exception to itself, cannot fabricate a human
    directive, and cannot suppress the exception silently.

  \item[$\bullet$] \textbf{Step 4: Pathway failure is a stall, not a collapse.}
    If $P(t) = \emptyset$ at time $t$, the \texttt{Escalate} algorithm cannot
    propagate the exception. It records the failure in the Forensic Ledger and the
    machine enters a stall condition — a non-terminal state — awaiting pathway
    restoration or manual intervention. No state transition leads to autonomous harm
    while $P(t) = \emptyset$.
\end{itemize}

Since every path to \text{RESOLVED} requires a human-issued directive (Step 1), since
entry to \text{ESCALATING} cannot be suppressed by any machine actor (Step 2), since
the escalation chain cannot terminate at a machine actor (Step 3), and since pathway
failure produces a stall rather than autonomous continuation (Step 4), it follows that
autonomous collapse — the propagation of an unresolved exception exclusively through
machine actors to a machine-terminal state — is structurally impossible in any
\bxthree{} system executing $\mathcal{B}$. $\square$

% ----------------------------------------------------------------
\subsection{LIVENESS: Eventual Human Contact}
\label{sec:liveness}

\textbf{Theorem (Liveness).} \emph{For any trigger event $e$ processed by
$\mathcal{B}(n)$, if at least one functional communication pathway to $h(n)$ exists at
some time $t \geq t_0$ (where $t_0$ is the trigger firing time), then the protocol
reaches the \text{HUMAN\_ACKNOWLEDGED} state within a bounded, finite wall-clock time,
or reaches \text{RESOLVED} via a timeout with an \texttt{unresolved} tag that records
the absence of human contact.}

\medskip
\noindent\textbf{Formal statement.} Let $t_0$ be the time at which T1 fires for node
$n$. Let $D_{\max}$ be the maximum depth of the parent chain from $n$ to $h(n)$. Let
$T_{\max}$ be defined as in Equation~\ref{eq:t-max}. Then:

\begin{equation}
\label{eq:liveness}
\exists\, t \in [t_0,\; t_0 + T_{\max}] :
\text{state}(n) \in \{\text{HUMAN\_ACKNOWLEDGED},\ \text{RESOLVED}\}
\;\land\;
\text{ledger}(t) \;\text{contains a human-anchored entry}.
\end{equation}

In all cases — successful human acknowledgment, explicit human timeout, or pathway
exhaustion — the liveness property is satisfied: the protocol reaches a definitive
state with a full forensic record, and no execution can stall indefinitely without
recording its stall as a protocol event.

\medskip
\noindent\textbf{Proof sketch.} The proof establishes a finite upper bound on the time
to human contact by analyzing each phase of the protocol sequentially.

\begin{itemize}
  \setlength\itemsep{0pt}
  \item[$\bullet$] \textbf{Step 1: Bounded time in BOUNDED.}
    T1 fires immediately upon trigger evaluation. The node has at most
    $T_{\text{bounds}}$ to produce a \fact-layer-approved resolution before T2 fires
    and transitions to \text{BAIL\_PENDING}. This is a hard timeout enforced by the
    \fact{} layer pre-commit hook — it cannot be suppressed or extended by any machine
    actor.

  \item[$\bullet$] \textbf{Step 2: Atomic transition to \text{ESCALATING}.}
    T4 ($\text{BAIL\_PENDING} \rightarrow \text{ESCALATING}$) fires immediately upon
    entry to \text{BAIL\_PENDING} with no dwell time. The wall-clock delta from T3
    firing to entry in \text{ESCALATING} is bounded by the atomic state transition
    latency — a deployment constant independent of node behavior.

  \item[$\bullet$] \textbf{Step 3: Bounded propagation along the parent chain.}
    In \text{ESCALATING}, the \texttt{Escalate} algorithm traverses the immutable parent
    chain upward. Each hop moves from node $m$ to $\alpha(m)$, strictly reducing
    graph-theoretic depth. Since the parent graph is a rooted tree with root
    $h_{\text{root}}$ and maximum provisioned depth $D_{\max}$, the number of hops to
    $h(n)$ is at most $D_{\max}$. For each hop, \texttt{SendToParent} retries up to
    $N_{\max}$ times with exponential backoff capped at $T_{\text{cap}}$. The total
    propagation time per hop is therefore bounded by $N_{\max} \cdot T_{\text{cap}}$.

  \item[$\bullet$] \textbf{Step 4: Alternate pathway fallback.}
    If a parent node does not acknowledge within $T_{\text{cap}}$, the algorithm fires a
    \texttt{parent\_unreachable} event and advances to the next available functional
    pathway in the Pathway Health Registry. This ensures the protocol cannot dead-end
    at a non-responsive ancestor. The fallback pathway selection is starvation-free: if
    any functional pathway to $h(n)$ exists, the algorithm will eventually select it.

  \item[$\bullet$] \textbf{Step 5: Human timeout as a liveness-preserving outcome.}
    Upon reaching $h(n)$, the protocol enters \text{HUMAN\_ACKNOWLEDGED} (T5). If the
    human anchor does not issue a directive before $T_{\text{human}}$ expires, the
    protocol transitions to \text{RESOLVED} with an \texttt{unresolved} tag (T8),
    recorded in the Forensic Ledger. Both outcomes — human acknowledgment and human
    timeout with \texttt{unresolved} tag — constitute human contact as defined by
    the liveness property.

  \item[$\bullet$] \textbf{Combined bound.}
    The total wall-clock time from $t_0$ to human contact is bounded by:
    \begin{equation}
    \label{eq:t-max}
    T_{\max} \;\leq\; T_{\text{bounds}} \;+\; D_{\max} \cdot N_{\max} \cdot T_{\text{cap}} \;+\; T_{\text{human}}.
    \end{equation}
    All terms are deployment constants. A deployment that requires a tighter bound
    must provision additional pathways, increase $N_{\max}$, or reduce $T_{\text{cap}}$
    and $T_{\text{human}}$ to bring the worst-case bound within the required limit.
\end{itemize}

Since each phase contributes a bounded, finite delay, and since the escalation chain
has bounded depth $D_{\max}$, the total time to reach a human-anchored state is bounded
by $T_{\max}$. No execution can loop indefinitely, retry indefinitely, or suppress the
timeout. The liveness property holds. $\square$

% ----------------------------------------------------------------
\subsection{ACCOUNTABILITY: Human Always Identifiable}
\label{sec:accountability}

\textbf{Theorem (Accountability).} \emph{For any directive $d$ recorded in the Forensic
Ledger as resolving a trigger event $e$, the issuer of $d$ is a human principal, and the
human principal's identity is uniquely and verifiably attributable in the ledger. No
machine actor can cause a directive to appear in the ledger under a falsified human
identifier.}

\medskip
\noindent\textbf{Formal statement.} Let $\mathcal{L}$ be the Forensic Ledger. Let
$e_{\text{resolve}}$ be the resolution event for trigger $\tau$. Then:

\begin{equation}
\label{eq:accountability}
\forall\, \tau,\; \forall\, d \in \mathcal{L}\ \text{recorded as resolving}\ \tau :
\text{issuer}(d) \in \text{HUMAN} \;\land\;
\text{ledger}^{-1}(d) = \tau \;\land\;
\text{SHA256-chain}(\mathcal{L}, d) = \text{valid}.
\end{equation}

That is: every directive that resolves a trigger is attributed to a human, the chain
of custody from trigger to directive is complete, and the SHA-256 digest chain confirms
that no ledger entry has been modified retroactively.

\medskip
\noindent\textbf{Proof sketch.} The proof proceeds by structural analysis of the
ledger-write authority and the chaining mechanism.

\begin{itemize}
  \setlength\itemsep{0pt}
  \item[$\bullet$] \textbf{Step 1: The \fact{} layer pre-commit hook is the sole
    write authority for directives.}
    By Assumption~\ref{asm:fact-layer}, the pre-commit hook rejects any ledger entry
    that would attribute a directive to a non-human actor. This rejection is
    architecturally enforced — not a policy check, not a configurable rule. A machine
    actor cannot produce a ledger entry bearing a human identifier because the write
    will be rejected at the \fact{} layer before it is committed.

  \item[$\bullet$] \textbf{Step 2: The \purpose{} layer carries the human anchor but
    cannot be bypassed by any machine actor.}
    The \purpose{} layer's role is exclusively receiving exceptions and issuing
    directives — it cannot authorize ledger entries attributing a directive to a
    non-human actor. This separation of concerns is enforced architecturally: the
    \purpose{} layer has write access to directive issuance, but the \fact{} layer
    pre-commit hook independently verifies that every ledger write's issuer field
    contains a human identifier before committing the entry.

  \item[$\bullet$] \textbf{Step 3: Atomic transition recording.}
    Every transition in $\mathcal{B}$ fires atomically with a corresponding ledger
    write: the \fact{} layer pre-commit hook commits both the state update and the
    ledger entry in a single atomic transaction. There is no state change without a
    corresponding ledger entry, and no ledger entry without a corresponding state
    change. Partial recording is architecturally impossible.

  \item[$\bullet$] \textbf{Step 4: Immutability of the human anchor.}
    By Assumption~\ref{asm:immutable-pointer}, $h(n)$ is set at node provisioning and is
    read-only at the machine-actor level for the node's operational lifetime. A machine
    actor cannot modify $h(n)$ to redirect accountability to a different human at
    runtime. Any attempted modification is rejected and logged as an unauthorized
    modification attempt in the Forensic Ledger. This prevents a compromised or
    malfunctioning node from substituting its own accountability anchor.

  \item[$\bullet$] \textbf{Step 5: SHA-256 chain integrity.}
    Every protocol event — trigger evaluations, state transitions, pathway selections,
    acknowledgments, directives, and timeout conditions — is committed to the Forensic
    Ledger as an append-only entry. Each entry is linked to its predecessor by a
    SHA-256 hash, forming a tamper-evident chain. Retroactive modification of any entry
    would break the hash chain and be computationally detectable. The chain of custody
    is complete: any auditor can reconstruct the full event sequence, attribute every
    state change to the principal who caused it, and confirm that no machine actor
    intervened between any node and its human anchor.
\end{itemize}

Accountability ensures attribution, not correctness: it guarantees that every directive
in the ledger is attributable to a human principal, but it does not guarantee that the
human principal made the correct decision. The latter is a separate concern — of
human factors, organizational process, and training — that the protocol cannot address
architecturally. What the protocol guarantees architecturally is that no machine actor
can issue, simulate, or substitute a directive without this being detectable as a
protocol violation. The ledger provides the documentary evidence; the pre-commit hook
provides the enforcement. Together they make accountability non-repudiable: a human
principal cannot later deny issuing a directive that appears in the ledger with their
identifier, and a machine actor cannot cause a directive to appear in the ledger under
a falsified human identifier. $\square$

% ----------------------------------------------------------------
\subsection{The Bailout Invariant}
\label{sec:bailout-invariant}

The three properties are not independent axioms — they form a logically coherent
constraint system. SAFETY establishes that the human anchor $h(n)$ exists and is
structurally reachable. LIVENESS establishes that $h(n)$ is reached within a bounded,
finite number of propagation steps. ACCOUNTABILITY establishes that $h(n)$ is uniquely
and verifiably identified in the Forensic Ledger. Their conjunction establishes the
\emph{Bailout Invariant}:

\beginInvariant}[Bailout Invariant]
\label{inv:bailout-invariant}
For any exception event $e$ processed by the Bailout Protocol in a \bxthree{} system,
there exists exactly one identifiable Human Accountability Anchor $h$ who bears
accountability for the outcome of $e$, and $h$ is reachable within a bounded, finite
number of propagation steps.
\endInvariant}

The Bailout Invariant holds unconditionally — regardless of the state of the machine
actor layer, the depth of the recursive hierarchy, or the presence of network
partitions. It is a structural guarantee, not a statistical one. Its satisfaction is
enforced by the \fact{} layer pre-commit hook on every state transition of
$\mathcal{B}$.

\medskip
\noindent\textbf{Classical safety-liveness mapping.} Classical distributed systems
theory distinguishes safety properties (``nothing bad ever happens'') from liveness
properties (``something good eventually happens''). The Bailout Protocol's three
properties map onto this framework as follows: SAFETY is a safety property — autonomous
collapse is a permanent, irreversible bad outcome that is architecturally impossible.
LIVENESS is a liveness property — human contact is a good outcome that must eventually
occur within a bounded time. ACCOUNTABILITY is neither purely safety nor purely
liveness — it is a \emph{knowledge property} ensuring that responsibility is not just
achieved but \emph{attributed}. It closes the loop between action and identification,
and it is this property that distinguishes the Bailout Protocol from conventional
exception-handling frameworks, which typically address only safety and liveness without
mechanisms for unique, cryptographically verifiable human identification in an
autonomous multi-agent hierarchy.

The Bailout Invariant is thus the conjunction of a safety property, a liveness
property, and a knowledge property — a combination that is rarely achieved in
distributed systems because the safety-liveness tradeoff (the FLP impossibility result)
typically forces a choice. The Bailout Protocol sidesteps the tradeoff by relying on a
strongly ordered, deterministic architecture with a fixed communication substrate, which
is precisely the setting in which all three property classes can hold simultaneously.
This is not a coincidence: it is a deliberate consequence of the \bxthree{} framework's
decision to encode accountability structurally rather than electorally.